\section{The fixed-point setup}
\label{sec:setup}

Every pool operation triggers a sync, and every sync must publish one
number: the senior share rate $y$. \Cref{sec:overview:dawn-to-dusk}
showed why this is not a plain oracle read. The rate depends on the
junior tranche's collateral value, the collateral contains senior
shares, and so the rate depends on itself. This section turns that
circularity into a single equation in one unknown and reduces the whole
problem to two facts about one function.

\subsection{What we must prove}
\label{sec:setup:roadmap}

Write the senior NAV the sync must produce as $x$. We will see that a
self-consistent $x$ is exactly a zero of the \emph{error function}
\[
  E(x) \;=\; F\!\left(P_A(x / N)\right) - x,
\]
where $P_A$ values the junior collateral at a candidate rate and $F$ is
the accountant's sync map. The on-chain solver finds that zero by
integer bisection, and bisection needs exactly two facts:

\begin{description}
  \item[(Monotonicity)] $E$ never rises: each step holds or falls. So the
    values where $E \geq 0$ form a prefix, the crossing is \emph{unique},
    and bisection knows which way to step.
  \item[(Bracketing)] $E$ is non-negative at the bottom of the search range
    and negative at the top. This makes the crossing \emph{exist} and
    brackets it.
\end{description}

The rest of this section defines the three objects $F$, $P_A$, and $E$.
\Cref{sec:existence} proves the two facts. \Cref{sec:solver} runs the
solver. Nothing else is load-bearing.

\subsection{The sync map $F$}
\label{sec:setup:F}

On each sync, the Dusk kernel calls the shared Royco accountant, whose
sync map we denote $F$. Given the
current raw NAVs of both tranches, $F$ runs the attribution and
settlement steps of \cref{sec:overview:waterfall} against the stored
checkpoint and returns the effective NAVs:
\[
  F\colon \bigl(\STraw,\ \JTraw;\ \text{checkpoint}\bigr) \longmapsto (\STeff,\ \JTeff).
\]
The two properties of $F$ that matter were stated in
\cref{sec:overview:waterfall}: it conserves NAV ($\STraw + \JTraw =
\STeff + \JTeff$), and it is piecewise-affine with slopes in $[0,1]$,
joined continuously at its kinks. NAV conservation holds by
construction, since every stored checkpoint is itself a prior output of
$F$. The monotone, 1-Lipschitz behaviour these two properties imply is
\cref{lem:F} in \cref{sec:existence}.

In a Dusk market the senior raw NAV is exogenous to the loop: it is the
value of the senior tranche's own assets, marked independently of the
pool. Only the junior raw NAV depends on the senior rate, because the
junior's collateral is the BPT. The senior side therefore drops out, and
for one sync, with $\STraw$ and the checkpoint fixed, we read $F$ as a
function of a single variable:
\[
  F\colon \JTraw \longmapsto \STeff.
\]

\subsection{The conservative valuation $P_A$}
\label{sec:setup:PA}

The junior tranche owns a fraction $f$ of a
Gyro E-CLP pool whose two legs are the senior share and a quote asset.
$P_A(y)$ is the value of that position when the senior share trades at
rate $y$, marked conservatively so that no manipulation and no rounding
can ever credit the junior tranche more than it could withdraw.

\paragraph{The frozen curve.}
The senior leg is priced by the rate the kernel publishes. At the start of a solve the
kernel reads the cached rate $y_0$ once and computes the pool invariant
$L$ from the rate-scaled balances, rounded down in favour of the senior
tranche. $L$ is then held fixed while the candidate rate $y$ varies.
Freezing $L$ is what makes $P_A$ a clean function of $y$ alone, which the
monotonicity proof relies on. The implementation details of the scaling
and rounding are collected in \cref{sec:solver:obligations}.

\paragraph{The fair point.}
Given $L$, the conservative value at candidate rate $y$ is the pool's
value at its \emph{fair point}: the point on the frozen curve whose
marginal price equals $y$. Write this value $V(y)$ and write
\[
  D(y) \;=\; V'(y),
\]
the pool's \emph{demand} for senior shares at price $y$, that is, the
senior reserve the fair point holds. Two geometric facts about the
E-CLP, proved in \cref{sec:existence:mono}, are all we need: $V$ is
non-decreasing, with $V'(y) = D(y) \geq 0$ (a higher senior price raises
the value of the senior shares the pool holds), and the demand $D(y)$ is
finite at every rate the curve can express.

\paragraph{The clamp.}
The fair point is a geometric construct, and far below the operating
rate the curve's demand $D(y)$ can grow without bound: the quote leg,
converted along the curve, would buy more senior shares than exist. That
is impossible. Even if every senior holder sold into the pool, it would
hold at most the total supply $N$, so the true demand never exceeds $N$.
We cap it there:
\[
  \Dc(y) \;=\; \min\bigl(D(y),\ N\bigr),
  \qquad
  \widehat{V}(y) \;=\; \int_0^y \Dc(u)\,du.
\]
By definition $\widehat{V}(0) = 0$, and $V(0) \geq 0$ since a reserve
is non-negative. With $\Dc \leq D = V'$ this gives
$0 \leq \widehat{V} \leq V$ everywhere: the cap only lowers the mark and
preserves conservatism. At any healthy
rate the live demand is far below $N$ and the cap is inactive. Its
purpose is to bound demand by the supply, which \cref{sec:existence}
turns into unconditional monotonicity.

\paragraph{The valuation.}
The junior tranche's pro-rata claim, rounded down in favour of the
senior tranche, is the \emph{conservative BPT valuation}:
\begin{equation}
  P_A(y) \;=\; \bigl\lfloor\, f \cdot \widehat{V}(y) \,\bigr\rfloor.
  \label{eq:PA}
\end{equation}
By the geometric facts and the clamp, $P_A$ is non-decreasing in $y$ and
Lipschitz with constant $f N$ (\cref{sec:existence:mono}). The value $V$
depends only on $L$ and $y$, never on the live pool composition, so a
swap cannot move it: a swap conserves $L$ in the fee-free pool math, and
the fee it leaves behind only grows $L$. Flash-loan rebalancing of the
pool cannot change the mark.

\subsection{The fixed-point equation}
\label{sec:setup:circularity}

The loop now closes in one line. To value the junior collateral the
kernel needs $\JTraw$, but $\JTraw = P_A(y)$, the rate $y = \STeff/N$,
and $\STeff = F(\JTraw)$. The senior NAV the sync must produce therefore
has to reproduce itself:
\begin{equation}
  x \;=\; F\!\left(P_A\!\left(x / N\right)\right).
  \label{eq:fixed-point}
\end{equation}
The loop is intrinsic, not an ordering bug: the junior's collateral
contains senior shares, so pricing it requires the senior price. Were
the junior's capital pure quote, $P_A$ would not depend on $y$ and
there would be no loop.

A solution of \eqref{eq:fixed-point} is a zero of the error function
$E$ (\cref{sec:setup:roadmap}), and the published rate is
$y^\star = x^\star / N$. \Cref{sec:existence} proves $E$ has
exactly one zero (Monotonicity) inside a range whose endpoints straddle
it (Bracketing).
